\section{Review Exercises}

\subsection{Elementary computational and graphical exercises}

\begin{enumerate}

\item For the circle
\[
(x-2)^2+(y+1)^2=16,
\]
find the centre and radius, sketch the curve, and write a trigonometric
parameterisation that traverses it once.

\item Rewrite
\[
x^2+y^2-6x+4y-3=0
\]
in standard circle form. Then describe the same circle in polar coordinates
centred at its own centre.

\item For the ellipse
\[
\frac{x^2}{25}+\frac{y^2}{9}=1,
\]
find $a$, $b$, $c$, the foci, and the eccentricity. Write the parameterisation
$x=5\cos t$, $y=3\sin t$ and plot several points.

\item Rewrite
\[
9x^2+4y^2-36x+8y+4=0
\]
in standard ellipse form. Find its centre, semi-axes, and eccentricity.

\item A parabola has focus $F=(2,0)$ and directrix $x=-2$. Derive its Cartesian
equation and write an algebraic parameterisation.

\item For the hyperbola
\[
\frac{x^2}{16}-\frac{y^2}{9}=1,
\]
find the vertices, foci, eccentricity, and asymptotes. Sketch both branches.

\item Convert the polar equation
\[
r=4
\]
to Cartesian form and describe the curve.

\item Convert
\[
r=6\cos\theta
\]
to Cartesian form. Identify the centre and radius and sketch the circle together
with the pole.

\item For
\[
r=\frac{6}{1+\frac12\cos\theta},
\]
identify $p$ and $e$, classify the conic, and compute the nearest and farthest
distances from the focus.

\item Classify each polar conic and state whether the trajectory is bounded:
\[
r=\frac{3}{1+0.4\cos\theta},
\qquad
r=\frac{3}{1+\cos\theta},
\qquad
r=\frac{3}{1+1.5\cos\theta}.
\]

\end{enumerate}

\subsection{Development, visualisation, and reasoning exercises}

\begin{enumerate}

\item Draw one diagram that compares Cartesian and polar coordinates for the same
point. Label $x$, $y$, $r$, and $\theta$, then derive
\[
x=r\cos\theta,
\qquad
y=r\sin\theta.
\]
Explain which coordinate system is better adapted to a distinguished focus.

\item Starting from the focus--directrix condition
\[
d(P,F)=e\,d(P,D),
\]
place the pole at $F$ and the directrix at $x=d$. Derive
\[
r=\frac{p}{1+e\cos\theta},
\qquad p=ed.
\]
Include a labelled geometric diagram showing the distance to the directrix.

\item Fix $p=2$ and sketch or plot the polar curves corresponding to
\[
e=0.2,
\qquad e=0.6,
\qquad e=0.9.
\]
Describe how the nearest and farthest focal distances change as eccentricity
increases.

\item For an ellipse with semi-major axis $a$ and eccentricity $e$, use
\[
p=a(1-e^2)
\]
to prove
\[
r_{\min}=a(1-e),
\qquad
r_{\max}=a(1+e).
\]
Then solve these equations for $a$ and $e$ in terms of $r_{\min}$ and
$r_{\max}$.

\item An observed orbit has nearest focal distance $4$ and farthest focal distance
$12$. Determine its semi-major axis, eccentricity, and focal parameter. Write its
polar equation with the nearest point at $\theta=0$.

\item Compare the two descriptions of the ellipse
\[
\frac{x^2}{a^2}+\frac{y^2}{b^2}=1
\]
and
\[
r=\frac{a(1-e^2)}{1+e\cos\theta}.
\]
Explain what becomes easy to read in the centre-based Cartesian description and
what becomes easy to read in the focus-based polar description.

\item The parameterisation
\[
x=a\cos t,
\qquad
y=b\sin t
\]
and the polar equation
\[
r=\frac{a(1-e^2)}{1+e\cos\theta}
\]
both describe an ellipse. Explain why $t$ is not generally the same angle as
$\theta$. Support the explanation with a labelled drawing.

\item Study the rotated polar equation
\[
r=\frac{p}{1+e\cos(\theta-\theta_0)}.
\]
Explain separately the roles of $p$, $e$, and $\theta_0$. Sketch two curves with
the same $p$ and $e$ but different $\theta_0$.

\item Complete squares and reconstruct the geometry of
\[
x^2+4y^2-6x+8y-3=0.
\]
Find the centre, semi-axes, foci, eccentricity, and the focus-based polar equation
after translating the Cartesian origin to the centre and orienting the polar
axis along the major axis.

\item Explain why the classification
\[
0\leq e<1,
\qquad e=1,
\qquad e>1
\]
is especially useful for mechanics. Your answer should distinguish bounded
motion, threshold escape, and unbounded passage, while making clear that the
physical derivation of the orbit equation has not yet been performed.

\end{enumerate}
